\section{Measure the two traversals}\label{blit-measurement}
Both tested implementations obey the same contract. The preclipped implementation
walks the intersection after calculating its edges. The alternative visits every
source cell, bounds-checking the corresponding destination coordinate. Both skip
source loads for invisible cells in these implementations; the contract also
permits reading clipped-away logical source cells.

\begin{table}[htbp]
\centering
\caption{Complete function work for the supplied descriptors and colors. Instruction counts include entry, restore, and return; they exclude caller setup, assembly, session creation, display rendering, and exit.}
\tagpdfsetup{table/header-rows={1}}
\begin{tabular}{rrrrrr}
\toprule
$x$ & $y$ & Preclipped & Per-cell & Source reads & Stores\\
\midrule
$-2$ & 1 & 227 & 446 & 12 & 8\\
0 & 0 & 308 & 604 & 20 & 13\\
6 & $-2$ & 126 & 332 & 4 & 2\\
$-1$ & 4 & 170 & 434 & 8 & 6\\
8 & 0 & 44 & 306 & 0 & 0\\
\bottomrule
\end{tabular}
\end{table}
The two central columns count native simulator instructions, not assembly source
lines or processor cycles. The final columns count source-read and destination-store
notices; those counts are equal for the two implementations on each listed scene.
Descriptor loads and private stack accesses are outside those final two columns.
The same-color store-eliding equivalent may produce different destination counts
while retaining the required opaque count.

For the default placement, removing repeated coordinate bounds checks reduces
native instruction work from 446 to 227. For an invisible placement, the preclipped
version exits after its geometric test, while the alternative still examines the
source coordinates. These are measurements for the supplied workloads, not a claim
about every image or elapsed execution time. A different transparency pattern,
unrolling strategy, or access width changes the work. Preserve the complete behavior
contract before comparing another implementation.
